\chapter{Abstract mathematics}

Mathematics, arguably, in a sense, can be said to be the field of \textit{abstraction}. It is fairly apparent that the human mind abstract the quantification for purposes of counting and managing, objectification hence forth to analysis, geometry for analysis of quantified shapes, and so on. Abstraction has already started from the root of mathematics, as for it to build an interpretation of the logical conclusion of the world, in either our framework of patterns recognition, structural realization, or integrally logical systems that is apparent of the physical world. Thereby, there are many ways to do abstract things, and hence require the field of abstract mathematics as a whole, though its definition, per chance, differs from now to the past, and toward the future from the present. 

There exists two kinds of abstraction one can make at the most fundamental level of mathematics, or so call as myself indict on such stance. One is the object-based, or cardinal abstraction, of which treating and examining between objects and their relations, interaction, laws, fundamental `shape' and size, and so on. Another one is to present a \textit{representation} \index{mathematical representation} on the object of interest itself, and so on, as to embed it in such representation for analysis. Both of them go interchangeably throughout the history of mathematical development. Said can be examined and inquired of from books like \cite{10.5555/1513920}, or philosophical analysis like \cite{Shapiro1997-SHAPOM-18,sep-philosophy-mathematics}. In essence, one inquires the idea of \textit{form}, of which bounds the combinatorial object to its discreteness. Analysis is again, simply the addition to such abstraction using a representation that fits certain structural criteria, explanation, basis, and so on. Though irrelevant of our own discussion, historically, the shift between mathematics has always been the argument toward that if mathematics is universal, or it is in certain sense stems from experience of the mathematicians. Toward such, the personal opinion of mine is to think of such as a hybrid, since while it is true that such mathematical explorations are unacquitted of human emotions and cloudy mind, it is ultimately the human himself ponders the mathematics. Hence, one can say that it abstracts from experience itself, though that does not mean one cannot imagine, or determine from such abstractness, another theory or another system of interesting mathematics, and so on. It just means that you cannot have an out-of-the-world mathematical system, since such is for the human mathematician that creates it would have to be out of the world itself. After all, the mathematician is still human, and any such imagination cannot be further from the physical analogies or based of such analogy for extension. Quantification for example, strips bare object of the limitation of manifestation ($n$ apples instead of only available apples you can have) so that to ponder of infinity. But infinity is a strange concept that no one can truly grasp - and that is the point. Toward the universal theory, then Gödel's incompleteness theorem indirectly, from which one can see, there also exists no complete theory in which every statement can be deducted simply of its form in the theory itself. Thus mathematics is not perfect, and hence one can see it as a fruitless endeavour to try forming one, for every loophole closed there are a few more in its wake. Though, that said, it also means mathematicians will still have their job security in the long run. And perhaps one might see satisfaction in the unsatisfactory outcome of such mathematical pursuit. 

And hence, after such, we can determine, in the barebone of the field, two interesting fields that in one way or another, represents such outlook. The first one is \textbf{abstract algebra}, of which concerns with objects, and, indeed, relations between them. And the other one is \textit{mathematical analysis} (there is also abstract geometry, analytical geometry and so on, but not that for now), and a few more, of which is concerned of such object, in a specific structural representation, on properties that is relevant of such treatment (convergence, continuity, etc.). While this chapter is about abstract mathematics, I would like to reserve it on the topic of abstract algebra itself, of which we retain the analysis part prominent to artificial intelligence later on in its own chapter for the reduction of concentrated confusion. And thus, we continue into exploring the world of abstract algebra. 

\section{The algebra of objects}

The set of object (yes we are talking about set theory then) $\{O\}$, is useless without any structure or relation on it. For example, an apple is usable because we can eat it, act on it, and so it can do what to us. Interaction and manipulation, structural description and so of are some of such way to describe and probe the particular object itself. One such way to do that, is called \textit{algebra}\index{algebra}. Per \cite{sep-algebra-logic-tradition}, algebra is the first abstraction one can get from the theory of arithmetic, of which then, manipulation of the object without making assumption of its form. Hence, algebra infers operations and thereof that examine \textit{non-destructive} manipulation of copula of objects in said set, for destructive means decomposition to lower components. Such is true, perhaps, from examples of such, and by extending toward various algebra in-between different object set itself. For example, there exists the $\mathbb{N}$-algebra, $\mathbb{Z}$-algebra, and so on up to $\mathbb{R}$-algebra of the real number set objects. It is also true, if one can see in historical and rhetorical sources of reference, like in \cite{sep-algebra-logic-tradition,Boole_1854,Russell_1919,Rashed_2009,Boole_Everest_1909} and so forth. One can also then observe the particular interesting fact that those operations in algebra is very intuitive, and familiar toward which one can directly take, for example, the 2-set operation $(\times, +)$ of a given algebraic structure as the addition and multiplication in reality, though much more abstract here. That is the spirit of quite literally, well, anything with algebra on objects. 

Historically, the origin of algebra is embedded into the historical stages in the evolution of mathematics. In such scheme, as indicted of \cite{Pinter2019ABO}, elementary algebra corresponds to the great classical age of algebra, which spans about 300 years from the sixteenth through the eighteenth centuries. It was during these years that the art of solving equations became highly developed and modern symbolism was invented. The world "Algebra" --- \textit{al jebr} in Arabic --- was first used by Mohammed of Kharizm, who taught mathematics in Baghdad during the ninth century. The word may be roughly translated as "reunion", and describes his method for collecting the terms of an equation in order to solve it. It is an amusing fact that the word "algebra" was first used in Europe in quite another context. In Spain, barbers were called \textit{algebristas}, or bonesetters (they \textit{reunited} the broken bones), because medieval barbers did bonesetting and bloodletting as a sideline to their usual business. The origin of the word clearly reflects the actual content of algebra at that time, for it was mainly concerned with ways of solving equations. In fact, Omar Khayyam, who is best remembered for his brilliant verses on wine, song, love and friendship which are collected in the \textit{Rubaiyat} - but who was also a great mathematician --- explicitly define algebra as the \textit{science of solving equations}. 

In line with such historical precedence, we might as well think of what constitute such popularity of manipulation, reunion of objects, and so on that algebra got, and why it is then called the science of equation solving. Surprisingly, a lot of foundational concept that can be shed light on with further rigours of the present time, can be found in the treatment in which algebra takes forth. In algebra, it then exists the concept of \textit{variable}, the concept of separating quantified language of numbers within such variable, and the concept of forming \textit{relation} on which algebra can then be settled. The language of number and encoding of such representations are populated of arithmetics, of which deal with them directly by using the language of manipulation and action operators. It is then required the abstraction of variable on top to work on manipulating such quantity, and another language to represent such abstractness. This is the basis for the quantities such as $x,y,z$ etc, and so on. Even numbers that are so often said as constant, are still denoted by $a,b,c$ and the non-arithmetics language, though their difference to variable $x$ is the scope in which they are applied to represent change (in an explicit form, $a,c,b$ are replaced with components from arithmetics, but $x,y,z$ remain the same, function as arbitrary global value holder). Lastly, the concept of relation, in which allows the expression like $x^{3}+ax^{2}+bx$ alongside the language of the variable abstraction together, forms the equation, of the $=$ equality relation, such that to relates different quantities together.

The discovery, evolution and schemes in which mathematical algebra was formed came forth during the Renaissance --- an age of high adventure and brilliant achievement, when the wide world was reawakening after the log austerity of the Middle Ages. Those men of intellectual pursuit who brought algebra to a high level of perfection at the beginning of its classical age, were as colourful and extraordinary a lot as have ever appeared in a chapter of history. Some of them, Cardan, Tartaglia (\textit{Ars Magna}), Ludovico Ferrari, and Niels Abel, brought the classical age to its end, and fruitfully provided the world with its wonder in the age of classical mathematics. It is not an exaggeration to pertain to the thought, that those men, arrogant and unscrupulous, brilliant, flamboyant, swaggering, and remarkable, lived their lives as they did their work: with style and panache, in brilliant dashes and inspired leaps of the imagination.

The modern landscape of algebra switched to a different direction, pondering different notion and thereof, after the beginning of the end that Niels Abel has concluded. The start of debates coming up switch to different variations of algebra, of which is indicted on the variations of \textit{representation} and \textit{relational object} itself, just as equality is, and transformation in \textit{linear algebra} might be. Those endeavours, discovering the new object of \textit{matrix}, the kind of Boolean algebra, and so on, open up to the old classical representation another peculiar, erratic and potentially remarkable direction to develop forward. Such is called unusual algebra, yet they exist on the basis of context, on the basis of axioms, system dynamics, and so on, of which lest the generality of classical notion, but with complexity never seen before. All of such development, requires the understanding and development of a unification of such modification -- of which then soon to be discovered, is the scheme of generalizing everything into the kind of arbitrary object that is called the \textbf{algebraic structure}\index{algebraic structure}. The main goal by then, is to say, to see and reveal the principle and framework on all known and all possible algebras, for which any further algebra can be created and understood, by variations and means of analysis that is available by existence. Under all of such evolution and advancement, also lies the hurdle of \textbf{abstraction}, in which now we strip of the old form of algebra as we know, and proceed on a greater generality in which we can focus. At best, abstraction is still something to focus on, and by itself, abstract has already begun in the simplest level as a way to quantify, so such is to say, it is a natural evolution that will inevitably be made, as far as any outlook is formed. The higher the abstraction, the more concrete the subject must be, and the more invariant and rigid the foundation should have. Such is where axioms are in play, thought I doubt talking more about it at this stage will provide any given insight into the problem onward. 

Most of our content here will be taken, for such requires solid foundation and the giant's shoulder for contents, books and materials like \cite{Dummit_Foote_2004,Artin_2011,Herstein_1975,Hungerford_1974,Gallian_2013,vanderWaerden_1930,Judson_2023,MacLane_1978}. Introductory undergraduate text in accordances are \cite{Cameron_2008,Armstrong_2010,Beardon_2005,Whitehead_2003,Carter_2009,Smith_Tabachnikova_2000}, and \cite{MathLibreTexts_Algebra}. It is wise to note that this chapter is the prerequisite before doing category theory, as one of the worry in tackling such is not the textbook or the concept is fairly self-contained, but because category theory is the abstraction level above, hence if one does not have experience on what is below, they cannot see why the abstraction is even worth it (usually). Preliminary topics should have been covered in the foundational chapter, and if not, would be present in this chapter in the appendix. 

\section{Operations}

A component that we did not mention that requires formalization too, in the sense of algebra, is the elementary operation set. Addition, subtraction, multiplication, division --- these and many others are familiar examples of operations on appropriate object, mainly sets of numbers. Intuitively, base on the definition, an operation on a set $A$ is a way of combining $n$-object from $A$ to form $m$ object on $A$. Thereby, we can regard them as internal transformation inside the set structure. This implies external transformation, and indeed, in the sense of function, a function that transform elements of $A$ to another set $B$ is regarded so, but that is not of our concern for now, as we talk about internal actions. 

Every operation is denoted by a symbol, such as $+,\times,\div$, and so on. Our purpose in this chapter is abstract, however --- we are not fixated on specific operation, but the generality of them all, and discover fact pertaining to operations generally rather than specific. Accordingly, we will show the power of such in the sense of creating our own operations, sometimes using those elementary operations that was listed, and act on arbitrary set of different interpretations. 

Let us begin by defining the notion of the operation. 
\begin{definition}[Operation]
    Let $A$ be any set of structured notion. Then, an \textit{operation}\index{algebraic operation} $*$ is a rule which assigns to each ordered pair $(a,b)$ of element of $A$ exactly one element $a*b$ in $A$. 
\end{definition}

The notation is fairly dummy, so in some texts, one use $aRb, a\cdot b, a\odot b$, and so on, but they are generally the same (use as your own preference). There are many rules which look deceptively like operations but are not, because this condition fails. Often $a*b$ is defined for all the obvious choice of $a$ and $b$, but remains undefined in a few exceptional cases. For example, division does not qualify as an operation on the set $\mathbb{R}$ of the real numbers, for there are ordered pairs such as $(3,0)$ whose quotient $3/0$ is undefined. In order to be an operation on $\mathbb{R}$, division would have to associate a real number $a/b$ with every ordered pair $(a,b)$ of element of $\mathbb{R}$. No exception allowed in such, and hence we said that $a*b$ must be \textbf{uniquely defined}. In other words, the value of $a*b$ must be given unambiguously. This means, again, result given without probabilistic approach (uncertain of such), not undefined, and only have one specific evaluation. Hence, the division operation is actually a kind called \textit{partial operation} \index{partial operation}, which hold such exception in hand. We then also said that
\begin{quote}
    If $a$ and $b$ are in $A$, $a*b$ must be in $A$. 
\end{quote}
This condition is often expressed by saying that $A$ is \textit{closed}\index{closure} under the operation $*$. This forbids the dilemma in which it breaks the internal restriction of the set, and turn into an external operation instead, and takes us out of $A$. 

An operation is \textit{any rule} which assigns to each ordered pair of elements of $A$ a unique element in $A$. Therefore, it is obvious that there are, in general, many possible operations on a given set $A$. Thereby, positional arguments are fairly natural for such pairing $(a,b)$. In total, we have 4 pairing possible for any arbitrary $a,b\in A$, and for example, the possible operation law: 
\begin{align}
    (a,a) = a ,\quad (a,b) = b , \quad (b,a) = a , \quad (b,b) = b \\
    (a,a) = a ,\quad (a,b) = a , \quad (b,a) = a , \quad (b,b) = a
\end{align}

And so on so forth. Each of the row describes a different operation on $A$. There then is said to exist 16 combinations possible for filling all of those operational rule, corresponding to the same action count. Continuing on, we can also apply some certain options for the rule on which the structure is determined, too. More so that unique properties that can be optional or not, but is general nonetheless. An operation, for example, may be \textit{commutative}\index{commutativity} if it can satisfy 
\begin{quote}
    An operation $*$ is commutative if $a*b=b*a$, for $a,b\in A$. 
\end{quote}
Sometimes, we might also can be \textit{associative}, that is, it may satisfy
\begin{quote}
    If $(a*b)*c = a*(b*c)$ for operation on $A$, and $a,b,c\in A$, then $*$ is associative. 
\end{quote}
To understand the importance of the associative law, it is basically a recursive expansion law. In a sense, operations are surely nice to be associative, because of their nested and recursive definition allowing us to treat pairing and nested operations to be exchangeable, hence making it easier to, as far as algebra is concerned, forcing terms out for rearrangement. If such associative law is enabled, it means constructing combination of three elements, without changing their order, in different applying order instead, and they produce the same result. If not, then we are left with no choice but to evaluate the nested operations on its own form. 

Those operations also highlight a fairly obscure and general, almost trivial fact that we choose to follow with the algebraic structure. That is, the \textit{type} of the operation. The condition of the operation $*: A\times A \to A$ itself indict that the object in combination, and the object in creation, is the same of the typing directive. That is, you would not be able to add an apple to an orange, and so on, in general wisdom. This is particularly valuable, and it also shows why defining generality like this would certainly help during the analysis and computational process, and something trivial of such is encoded straight into the definition on such generality. In cases of which operations indeed can add an apple to an orange, we can simply move the domain of computation to $C\supset A\times B$, and so on, which is borderline convenient. Unconsciously so, it is fairly similar to \textbf{type theory}, category theoretic approach, and the methodological principle of representation. Closure itself is the result of such choice, and hence the restriction on which triviality is stated of the exploratory operation. One then ask if this is inherently open, and can we do anything else about it, and the answer is yes, closure is not totally required. Nevertheless, there exists the fundamental notion of \textit{operational boundary}, in the sense of bounding actions to area of respective availability. Such external operation can be encoded in modern function concepts, and operation can take role to be the internal set transformation accordingly. It is wise to remember such notion. Though, jargon and technical terms might overlap, as we might see in the theory of \textit{operator}. 

Continuing, if there is an element $e\in A$ with the property that 
\begin{quote}
    $e*a=a$ and $a*e=a$ for every element $a$ in $A$. 
\end{quote}
then $e$ is called an \textit{identity}\index{identity} or "neutral" element with respect to the operation $*$. Roughly speaking, the equation tells us that when $e$ is combined with any element $a$, it does not change $a$, hence create a special \textit{invariant element} of set $A$. They differ from time to time in different operation, for example $e$ being $1$ in multiplication and $0$ in addition. Interestingly, if $a$ is any element of $A$, and $x$ is an element of $A$ such that
\begin{quote}
    $a*x = e$ and $x*a = e$
\end{quote}
then $x$ is called an \textit{inverse}\index{inverse} of $a$. Roughly speaking, it is a resetting element, that when an element is combined with its inverse, it produces the neutral element. Since neutral element is fairly elementary, any kind of such inverse will be the reverse operation toward such focus point. The inverse if often denoted as $a^{-1}$. 

\section{Algebraic structures}


